\section{Methodology}
\label{sec:method}

\subsection{Design principles}

The method was developed around four principles. First, generation is an
intermediate operation rather than the final product. Second, every technical
decision is made against a versioned asset contract. Third, feasibility and
preference are separated during selection. Fourth, human approval is required at
points where automatic metrics cannot adequately represent intent or aesthetic
judgment.

\Cref{fig:architecture} summarizes the resulting pipeline. A thin client calls the
same Agent runtime used by the command-line interface. The runtime coordinates
planning, generation, asset processing, rendering, evaluation, and persistence;
the user interface does not reimplement those operations.

\begin{figure}[t]
  \centering
  \begin{tikzpicture}[
    node distance=6mm and 8mm,
    box/.style={draw=skinsmithdark, rounded corners, align=center, minimum height=9mm,
      text width=27mm, fill=skinsmithlight, font=\small},
    checkpoint/.style={box, fill=skinsmithblue!12, draw=skinsmithblue},
    process/.style={box, fill=skinsmithteal!10, draw=skinsmithteal},
    arrow/.style={-{Latex[length=2mm]}, thick, draw=skinsmithdark}
  ]
    \node[box] (input) {Weapon + short theme};
    \node[process, right=of input] (theme) {Theme and style compilation};
    \node[checkpoint, right=of theme] (c1) {Checkpoint 1\\confirm theme};
    \node[process, below=of c1] (dirs) {3--4 textual directions};
    \node[checkpoint, left=of dirs] (c2) {Checkpoint 2\\select direction};
    \node[process, left=of c2] (art) {3--4 landscape artworks};
    \node[process, below=of art] (preview) {Low-cost UV mapping and three views};
    \node[checkpoint, right=of preview] (c3) {Checkpoint 3\\select mapped artwork};
    \node[process, right=of c3] (formal) {Formal weapon-space bake};
    \node[process, below=of formal] (eval) {Constraints, Pareto selection, Route C};
    \node[box, left=of eval] (export) {2048 PNG/TGA + evidence log};

    \draw[arrow] (input) -- (theme);
    \draw[arrow] (theme) -- (c1);
    \draw[arrow] (c1) -- (dirs);
    \draw[arrow] (dirs) -- (c2);
    \draw[arrow] (c2) -- (art);
    \draw[arrow] (art) -- (preview);
    \draw[arrow] (preview) -- (c3);
    \draw[arrow] (c3) -- (formal);
    \draw[arrow] (formal) -- (eval);
    \draw[arrow] (eval) -- (export);
  \end{tikzpicture}
  \caption{The \skinsmith\ workflow. The three blue nodes are blocking human
  checkpoints; formal processing cannot begin before a mapped artwork is selected.}
  \label{fig:architecture}
\end{figure}

\subsection{Agent state and human checkpoints}

An \texttt{AgentRunRequest} contains the asset identifier, a short theme or free
brief, an optional broad style family, the candidate count, and explicit budgets
for image calls, retries, and refinement rounds. The persisted state follows:

\begin{align*}
\texttt{created} &\rightarrow \texttt{planning}
\rightarrow \texttt{awaiting\_theme}
\rightarrow \texttt{awaiting\_direction}\\
&\rightarrow \texttt{awaiting\_artwork}
\rightarrow \texttt{executing}
\rightarrow \texttt{completed}.
\end{align*}

At the first checkpoint, a \texttt{ThemeCompiler} expands a compact keyword into a
validated \texttt{ThemePack}. The pack contains the concept, palette, primary and
supporting elements, environmental and material cues, atmosphere, component
story, evaluation criteria, and prohibited content. The user may confirm or edit
this interpretation. A \texttt{StyleCompiler} then produces three or four
materially different textual directions for the confirmed theme and selected
weapon. At the second checkpoint, the user chooses one direction.

The selected direction generates three or four dense landscape master artworks.
Each is validated for resolution, tonal and colour variation, local detail
density, and forbidden content. It is then passed through a low-cost mapping stage
using the target adapter. The original image and the left, right, and top views are
treated as one indivisible candidate card. The third checkpoint asks the user to
select one card. This design prevents a visually attractive source from being
chosen without seeing how it behaves on the asset.

The selected source path and hash become part of a frozen design contract. Formal
processing reuses that exact file; it does not regenerate an approximation after
selection. Every transition, tool action, budget change, retry, and stop reason is
written to persisted state so a failed or interrupted run can resume from the last
valid checkpoint.

\subsection{Asset adapter}

Each target is represented by a versioned \texttt{GameAssetAdapter}. The adapter
binds:

\begin{itemize}
  \item the mesh path, version, and SHA-256 hash;
  \item the authoritative UV source and UV address mode;
  \item canonical longitudinal and up axes;
  \item semantic regions and creative anchors;
  \item paintable-area rules;
  \item fixed cameras;
  \item texture size and export format.
\end{itemize}

This boundary is important because mesh geometry, UV charts, and material behavior
cannot be assumed to transfer between weapons. The accepted AK-47 adapter uses the
new CS2 geometry, a 2048 by 2048 RGB texture, fixed left/right/top cameras, and a
24-bit TGA Custom Paint export. The M4A4 case supplies a different mesh hash, axes,
UV topology, and seam calibration while reusing the core algorithms.

\subsection{Mesh-derived geometry maps}

Let a mesh triangle have three-dimensional vertices
$\mathbf{p}_0,\mathbf{p}_1,\mathbf{p}_2$, vertex normals
$\mathbf{n}_0,\mathbf{n}_1,\mathbf{n}_2$, and UV coordinates
$\mathbf{u}_0,\mathbf{u}_1,\mathbf{u}_2$. For every covered atlas pixel
$\mathbf{u}$, barycentric coordinates $(\lambda_0,\lambda_1,\lambda_2)$ are
computed in UV space. The corresponding position and normal are

\begin{align}
  \mathbf{p}(\mathbf{u}) &=
  \lambda_0\mathbf{p}_0+\lambda_1\mathbf{p}_1+\lambda_2\mathbf{p}_2,\\
  \tilde{\mathbf{n}}(\mathbf{u}) &=
  \lambda_0\mathbf{n}_0+\lambda_1\mathbf{n}_1+\lambda_2\mathbf{n}_2,\\
  \mathbf{n}(\mathbf{u}) &= \frac{\tilde{\mathbf{n}}(\mathbf{u})}
  {\lVert\tilde{\mathbf{n}}(\mathbf{u})\rVert_2}.
\end{align}

The mesh is transformed into a canonical frame whose $x$ axis follows the weapon
length, $y$ is up, and $z$ distinguishes the two sides. Positions are normalized
to the asset bounds. Rasterizing all UV triangles produces a position map, a
normal map, a valid-pixel mask, and a face-index map. These maps connect continuous
weapon-space composition to the existing fragmented atlas.

\subsection{Route A: tileable-pattern baseline}

\routeA\ generates a square, dense, crop-tolerant pattern. Its opposite image
borders are repaired using periodic-plus-smooth decomposition
\citep{moisan2011periodic}. The repaired square is then sampled through the
original UVs and rendered on the weapon. This route is deliberately generic: it
is expected to tolerate arbitrary crops, but it has no explicit whole-weapon
hierarchy.

Two seam concepts are recorded. The square-border error measures differences
between opposite image edges. The asset-seam error instead evaluates pairs of UV
edges that correspond to the same welded three-dimensional edge. Passing the
first measure does not imply passing the second.

\subsection{Route B: weapon-space composition and UV bake}

\routeB\ uses one selected landscape master artwork. The source is fitted to
continuous left, right, and top canvases in canonical weapon coordinates. The
source must be dense and edge-to-edge, with several medium or small thematic
clusters, because a single dominant subject or large empty region is unlikely to
survive the narrow silhouette and UV fragmentation.

For each valid UV pixel, the canonical position determines the sampling coordinate
on the side and top canvases. Surface normals blend the projections:

\begin{align}
  w_s &= |n_z|^p,\qquad w_t = \max(n_y,0)^p,\\
  \mathbf{c}(\mathbf{u}) &=
  \frac{w_s\mathbf{c}_s(\mathbf{p})+w_t\mathbf{c}_t(\mathbf{p})}
  {w_s+w_t+\epsilon},
\end{align}

where $p=4$ in the accepted configuration. $\mathbf{c}_s$ is sampled from the
appropriate left or right canvas according to surface orientation, while
$\mathbf{c}_t$ is sampled from the top canvas. This soft blend avoids an abrupt
switch between charts selected solely by the largest normal component.

\begin{figure}[t]
  \centering
  \includegraphics[width=\textwidth]{weapon_space_concept.png}
  \caption{The central Route B idea: compose in continuous weapon space, store
  fragments in the authoritative UV atlas, and judge coherence after
  reconstruction. The flat atlas is not expected to read as a conventional
  illustration.}
  \label{fig:weapon-space}
\end{figure}

\subsection{True asset-seam metric and correction}

OBJ files often duplicate vertices at UV boundaries. The method first welds
coordinate-identical geometric vertices, identifies mesh edges shared by adjacent
faces, and records pairs whose UV coordinates differ. These pairs form the true
asset-seam graph.

For $N$ paired samples with RGB colours $\mathbf{a}_i,\mathbf{b}_i$ and local
gradients $\mathbf{g}^{a}_i,\mathbf{g}^{b}_i$, the colour and gradient terms are

\begin{align}
  E_c &= \frac{1}{N}\sum_{i=1}^{N}
  \frac{\lVert\mathbf{a}_i-\mathbf{b}_i\rVert_1}{3\cdot255},\\
  E_g &= \frac{1}{N}\sum_{i=1}^{N}
  \frac{\lVert\mathbf{g}^{a}_i-\mathbf{g}^{b}_i\rVert_1}{3\cdot255}.
\end{align}

The implementation combines these normalized terms into $\assetseam$ using the
locked project weights. The principal hard requirement is
$\assetseam\leq0.01$.

An edge-safety correction blends a narrow region around UV-island boundaries
toward a stable base colour. Wider regions generally reduce seam error but can
erase visible detail. Candidate widths are therefore evaluated against both the
seam threshold and a minimum multi-view retention of 95\% relative to the
uncorrected texture.

\subsection{Rendering and evaluation}

The software renderer loads textured OBJ geometry, triangulates polygon faces,
interpolates UV coordinates barycentrically, samples the texture bilinearly, and
uses a z-buffer with fixed orthographic left, right, and top cameras. Lighting and
cameras are deterministic so candidates are compared under the same conditions.

Texture-level evidence includes contrast, saturation, detail, paintable coverage,
and square-border periodicity. The multi-view score combines foreground coverage,
contrast, saturation, visible detail, and luminance consistency across the three
views. The project also records mapped-element readability, defined as

\begin{equation}
  S_{\mathrm{read}} =
  0.25S_{\mathrm{source}}+
  0.50S_{\mathrm{left}}+
  0.15S_{\mathrm{right}}+
  0.10S_{\mathrm{top}}.
\end{equation}

This weighting reflects the principal inspection view in the case study. It is
recommendation-only and cannot override the user's artwork selection.

\subsection{Constraint-first selection}

For a candidate $x$, each hard constraint $j$ has a non-negative violation
$v_j(x)$. A candidate is feasible only if

\[
  V(x)=\sum_j v_j(x)=0.
\]

Feasible candidates are compared using Pareto dominance over the locked
objectives. If several candidates remain non-dominated, a deterministic priority
order selects the operating point. A conventional normalized weighted score is
retained only as an ablation baseline. This policy prevents a strong visual score
from compensating for a deployment failure.

\subsection{Route C: bounded local correction}

\routeC\ begins from the selected \routeB\ texture. Rendered component visibility
masks identify a component whose detail level deviates from its target. The Agent
creates exactly two correction candidates at fixed intensities. Changes are
restricted to that component plus a transition halo, and the number of changed
pixels outside the allowed region must be zero.

Each correction must retain all \routeB\ hard constraints. Let $S_B$ be the locked
score of the original and $S_C$ the best feasible correction. The correction is
accepted only when

\[
  S_C-S_B\geq0.01.
\]

Otherwise the final result is the unchanged \routeB\ texture. The method therefore
tests diagnosis and action without assuming that automatic refinement is
monotonic.

\subsection{Generation backends and evidence}

The generator interface supports a deterministic procedural baseline, a local
SD-Turbo model, and provider-backed image generation. SD-Turbo is pinned to
revision \path{b261bac6fd2cf515557d5d0707481eafa0485ec2} at 512 pixels, two
inference steps, and guidance scale zero. The accepted creative run used
\texttt{gemini-3.1-flash-lite} for structured planning and
\texttt{gemini-3.1-flash-image} for 1K landscape source images.

Every formal run preserves the request, prompts, provider and model identifiers,
seeds where applicable, redacted traces, source hashes, validation, mapped views,
UV outputs, metrics, decisions, budgets, retries, runtime, and checkpoints. API
keys are read from environment variables and are not written to run evidence.
